% A chronological reading list for game-based security proofs.
%
% Companion to files/all-uc-papers.pdf, which does the same thing for
% universally-composable security: there, each entry's sub-bullets name the
% ideal functionalities (F_com, F_zk, ...) the paper introduces or uses; here,
% they name the games, lemmas and transformations it introduces.
%
% Selection rule: representative papers that introduce a game-based security
% definition *formally* and prove something *formally* against it, chosen so
% that reading them in order shows how the style evolved. Not exhaustive, and
% not meant to be.
%
% Bibliographic data lives in refs.bib; every paper is \cite'd, so the body
% stays short and the References section carries the venues, DOIs and ePrint
% links. Build with `make` (see the Makefile in this directory).
\documentclass[11pt]{article}

\usepackage[T1]{fontenc}
\usepackage[utf8]{inputenc}
\usepackage{amsmath}   % \lvert, \rvert in the Difference/Fundamental Lemma statements
\usepackage{newtxtext}
\usepackage[margin=1in]{geometry}
\usepackage[dvipsnames]{xcolor}
\usepackage{enumitem}
\usepackage{titlesec}
\usepackage{microtype}
\usepackage[colorlinks=true,urlcolor=RoyalBlue,linkcolor=RoyalBlue,
            citecolor=RoyalBlue]{hyperref}
\usepackage{doi}       % turns the doi= field of a bib entry into a link

% Year headings, in the same washed-out grey as the "Crypto 2001" headings in
% all-uc-papers.pdf.
\definecolor{yearcolor}{gray}{0.40}
\titleformat{\section}{\Large\color{yearcolor}}{}{0pt}{}
\titlespacing*{\section}{0pt}{1.3\baselineskip}{0.4\baselineskip}

% Papers are numbered 1., 2., ...; what each one contributes is lettered
% a., b., ...; sub-points under those are i., ii., ... — same nesting as the UC
% list.
\setlist[enumerate,1]{label=\arabic*., leftmargin=2em, itemsep=2pt, topsep=3pt}
\setlist[enumerate,2]{label=\alph*., leftmargin=1.8em, itemsep=1pt, topsep=2pt}
\setlist[enumerate,3]{label=\roman*., leftmargin=1.8em, itemsep=1pt, topsep=2pt}

% What each paper contributes. A paper can carry more than one tag.
\newcommand{\kind}[1]{{\footnotesize\textsc{\textcolor{yearcolor}{[#1]}}}\ }
\newcommand{\defn}{\kind{defn}}
\newcommand{\method}{\kind{method}}
\newcommand{\tool}{\kind{tool}}
\newcommand{\turn}{\kind{turn}}
\newcommand{\case}{\kind{case}}

% Each numbered entry opens with the authors, the (short) title and the
% citation; the venue, year, DOI/ePrint link all come from refs.bib via the
% References section at the end.
\newcommand{\paper}[3]{#1, \emph{#2}~\cite{#3}.}

\title{\vspace{-2.5em}Game-Based Security Proofs\\[.2em]
       \large A chronological reading list}
\author{}
\date{}

\begin{document}
\maketitle
\vspace{-3.5em}

\noindent
Companion to \texttt{all-uc-papers.pdf}. That list walks the
universally-composable literature and, under each paper, names the ideal
functionalities it introduces or uses. This one walks the \emph{game-based}
literature and, under each paper, names the games, lemmas and transformations
it introduces.

\medskip\noindent
These are \emph{representative} papers, not all of them. A paper earns a place
here if it states a security notion as a formal game and proves something
formally about it --- and if reading it in sequence with the others shows how
the style got from one place to the next. It starts in 1982, with the first
security definition written as an experiment between an adversary and a
challenger, and ends with proofs of deployed protocols that no single game
sequence could hold. Full bibliographic data, with DOI and ePrint links, is in
the References at the end.

\medskip\noindent
Each entry is tagged:

\begin{itemize}[leftmargin=2em, itemsep=1pt, topsep=3pt]
  \item \defn states a new security notion as a formal game;
  \item \method a proof technique or framework for chaining games;
  \item \tool mechanizes that methodology (computer-aided / machine-checked);
  \item \turn a turning point --- a proof that broke or a definition shown
    inadequate, which changed how games were written afterwards;
  \item \case a landmark formal proof of a real primitive or protocol.
\end{itemize}

\medskip\noindent
\textbf{If you read only five.} \cite{GM82} (where the experiment comes from);
\cite{BDJR97} (what a fully formal, concrete game-based treatment of one
primitive looks like); \cite{Shou04} and \cite{BR04} (the method itself, in
prose and in a formal system); and \cite{SSP18} (how it is done at protocol
scale today).

\section{1982 --- the experiment appears}

\begin{enumerate}
  \item \defn \paper{Goldwasser \& Micali}{Probabilistic Encryption \& How to
    Play Mental Poker Keeping Secret All Partial Information}{GM82,GM84}
    \begin{enumerate}
      \item \textbf{Polynomial security}: the adversary names two messages,
        receives an encryption of one of them, and must say which --- the first
        security definition phrased as an \emph{experiment} between a
        challenger and an adversary. Renamed IND-CPA later.
      \item \textbf{Semantic security}: whatever is efficiently computable
        about the plaintext from the ciphertext is computable without it.
      \item The proof that the two are equivalent is arguably the first
        \textbf{game hop} in the literature: a security experiment rewritten
        step by step into a different one.
      \item The \textbf{hybrid argument}, used to go from single-bit to
        multi-bit messages. Every later multi-step game chain is this argument
        generalized.
    \end{enumerate}

  \item \defn \paper{Yao}{Theory and Application of Trapdoor Functions}{Yao82}
    \begin{enumerate}
      \item \textbf{Computational indistinguishability} of distribution
        ensembles: the relation that every ``these two games are close'' step
        is ultimately an instance of.
      \item Next-bit unpredictability $\Longleftrightarrow$ indistinguishability
        (with \cite{BM84}) --- an early proof that two differently-shaped games
        define the same thing.
    \end{enumerate}

  \item \defn \paper{Blum \& Micali}{How to Generate Cryptographically Strong
    Sequences of Pseudorandom Bits}{BM84}
    \begin{enumerate}
      \item The \textbf{next-bit test}: a \emph{prediction} game rather than a
        distinguishing game. The two shapes --- guess a hidden bit vs.\ produce
        a winning object --- are still the only two shapes in use.
    \end{enumerate}
\end{enumerate}

\section{1984 --- oracles enter the game}

\begin{enumerate}
  \item \defn \paper{Goldreich, Goldwasser \& Micali}{How to Construct Random
    Functions}{GGM84}
    \begin{enumerate}
      \item \textbf{PRF game}: distinguish oracle access to $F_k(\cdot)$ from
        oracle access to a uniformly random function. The first widely-used
        game in which the adversary \emph{queries an oracle adaptively} rather
        than receiving a fixed challenge.
      \item The GGM tree's analysis is the canonical multi-hop chain: one hybrid
        per level, each neighboring pair indistinguishable by the PRG game.
    \end{enumerate}

  \item \defn \paper{Goldwasser, Micali \& Rivest}{A Digital Signature Scheme
    Secure Against Adaptive Chosen-Message Attacks}{GMR88}
    \begin{enumerate}
      \item \textbf{EUF-CMA}: the adversary gets a signing oracle and wins by
        forging on any message it did not query. The template for every
        ``win condition + oracle'' game since.
      \item The \textbf{attack/break taxonomy} --- key-only, known-message,
        generic and adaptive chosen-message attacks, crossed with total break,
        universal/selective/existential forgery --- the first systematic grid
        of games for a single primitive.
    \end{enumerate}
\end{enumerate}

\section{1985 --- the road forks}

\begin{enumerate}
  \item \turn \paper{Goldwasser, Micali \& Rackoff}{The Knowledge Complexity of
    Interactive Proof Systems}{GMR89}
    \begin{enumerate}
      \item Zero knowledge is defined by exhibiting a \textbf{simulator}, not by
        bounding an adversary's advantage in a game. This is the branch that
        leads to UC~\cite{Can01} and to \texttt{all-uc-papers.pdf}.
      \item Included as the fork in the road: read next to \cite{GM82} to watch
        the two proof styles separate. Everything below is the game branch.
    \end{enumerate}
\end{enumerate}

\section{1988}

\begin{enumerate}
  \item \defn\method \paper{Luby \& Rackoff}{How to Construct Pseudorandom
    Permutations from Pseudorandom Functions}{LR88}
    \begin{enumerate}
      \item \textbf{PRP game} and \textbf{strong-PRP game} (the latter also
        gives the adversary the inverse oracle).
      \item The Feistel analysis proceeds through a sequence of idealized
        experiments, bounding the probability of a collision event that makes
        them differ --- game playing with a bad event, about fifteen years
        before it was named in \cite{BR04}.
    \end{enumerate}
\end{enumerate}

\section{1990--1991 --- the CCA ladder}

\begin{enumerate}
  \item \defn \paper{Naor \& Yung}{Public-Key Cryptosystems Provably Secure
    Against Chosen Ciphertext Attacks}{NY90}
    \begin{enumerate}
      \item \textbf{IND-CCA1} (``lunchtime''): a decryption oracle, available
        only \emph{before} the challenge ciphertext is handed over.
    \end{enumerate}

  \item \defn \paper{Rackoff \& Simon}{Non-Interactive Zero-Knowledge Proof of
    Knowledge and Chosen Ciphertext Attack}{RS91}
    \begin{enumerate}
      \item \textbf{IND-CCA2}: the decryption oracle stays available after the
        challenge, with the challenge ciphertext itself excluded. That exclusion
        clause is the single most-revisited line in any security game ---
        see~\cite{BHK09}.
    \end{enumerate}

  \item \defn \paper{Dolev, Dwork \& Naor}{Non-Malleable Cryptography}{DDN91}
    \begin{enumerate}
      \item \textbf{NM-CPA / NM-CCA}, defined by comparing the real adversary
        against a simulator that never sees the ciphertext --- a definition
        deliberately straddling the game and simulation styles, later also
        given a purely game-based (comparison-free) form.
    \end{enumerate}
\end{enumerate}

\section{1993 --- concrete security, and the first protocol games}

\begin{enumerate}
  \item \defn\method \paper{Bellare \& Rogaway}{Random Oracles are Practical: A
    Paradigm for Designing Efficient Protocols}{BR93rom}
    \begin{enumerate}
      \item The \textbf{random oracle model}: the hash function becomes an
        oracle \emph{inside} the game, which the reduction may observe and
        program. Hugely enabling --- and the source of most later disputes
        about what a game-based theorem means (\cite{CGH04}, \cite{BHK13}).
      \item \textbf{Concrete security}: theorems stated as explicit
        $(t, q, \epsilon)$ relations rather than asymptotics, which is what
        makes a game's advantage function the deliverable.
    \end{enumerate}

  \item \defn \paper{Bellare \& Rogaway}{Entity Authentication and Key
    Distribution}{BR93ke}
    \begin{enumerate}
      \item The \textbf{BR key-exchange game}: sessions modeled as oracles
        $\Pi_{i,j}^s$, with \textbf{Send}, \textbf{Reveal} and \textbf{Corrupt}
        queries, plus a single \textbf{Test} query answered with either the real
        session key or a random one.
      \item \textbf{Matching conversations} as the partnering predicate --- the
        bookkeeping that says which sessions are supposed to share a key.
      \item Directly ancestral to \cite{CK01}, \cite{JKSS12} and \cite{FG14}.
        If you want to see where protocol security games come from, this is the
        paper.
    \end{enumerate}
\end{enumerate}

\section{1994--1997 --- the concrete-security school}

\begin{enumerate}
  \item \method\case \paper{Bellare, Kilian \& Rogaway}{The Security of Cipher
    Block Chaining}{BKR94}
    \begin{enumerate}
      \item PRF-security of CBC-MAC with an explicit birthday bound
        ($\approx q^2\ell^2/2^n$): the first influential proof whose \emph{point}
        is the shape of the advantage function, not merely that it is negligible.
    \end{enumerate}

  \item \defn\turn \paper{Bellare \& Rogaway}{Optimal Asymmetric Encryption}
    {BR94}
    \begin{enumerate}
      \item \textbf{Plaintext awareness}: an adversary that produces a valid
        ciphertext must ``know'' the plaintext, formalized via an extractor
        reading its random-oracle queries.
      \item The claimed IND-CCA2 implication is the famous one that did not
        hold --- see \cite{Shou01}.
    \end{enumerate}

  \item \defn\case \paper{Bellare, Canetti \& Krawczyk}{Keying Hash Functions
    for Message Authentication}{BCK96}
    \begin{enumerate}
      \item MAC security as a formal game, and the reduction of HMAC to two
        assumptions on the compression function (PRF-ness, weak collision
        resistance) --- rather than to an unanalyzed ``hash is good'' wish.
      \item The paper that made the whole concrete-security style credible to
        practitioners, because the object it analyzed shipped everywhere.
    \end{enumerate}

  \item \method \paper{Kilian \& Rogaway}{How to Protect DES Against Exhaustive
    Key Search}{KR96}
    \begin{enumerate}
      \item A proof built as a sequence of games that are \textbf{identical
        until a flag is set}, with the whole difference charged to that flag.
        \cite{BR04} points at this paper as where their framework came from.
    \end{enumerate}

  \item \method \paper{Pointcheval \& Stern}{Security Proofs for Signature
    Schemes}{PS96}
    \begin{enumerate}
      \item The \textbf{forking lemma}: run the adversary twice on the same
        random tape with divergent oracle answers. Rewinding as a reusable game
        transformation; generalized in \cite{BN06}.
      \item The \textbf{splitting lemma}, the probability bookkeeping that makes
        the forking argument go through.
    \end{enumerate}

  \item \method \paper{Bellare \& Rogaway}{The Exact Security of Digital
    Signatures --- How to Sign with RSA and Rabin}{BR96}
    \begin{enumerate}
      \item \textbf{Tightness} promoted to a design goal: the scheme is chosen
        so that the game hop loses as little as possible. Sharpened
        in~\cite{Cor00}, bounded from below in \cite{BJLS16}.
    \end{enumerate}

  \item \defn \paper{Bellare, Desai, Jokipii \& Rogaway}{A Concrete Security
    Treatment of Symmetric Encryption}{BDJR97}
    \begin{enumerate}
      \item Four formalizations of chosen-plaintext security ---
        \textbf{find-then-guess}, \textbf{left-or-right},
        \textbf{real-or-random} and \textbf{semantic} --- each a precisely
        written game, with the reductions between them and their exact advantage
        losses.
      \item Concrete PRF-based analyses of ECB, CBC, CTR and friends in one
        common framework.
      \item The single best specimen of what a complete, formal, game-based
        treatment of one primitive looks like; still the template modern papers
        copy.
    \end{enumerate}

  \item \defn\method \paper{Shoup}{Lower Bounds for Discrete Logarithms and
    Related Problems}{Shou97}
    \begin{enumerate}
      \item The \textbf{generic group model}: group elements become opaque
        handles served by an oracle, so an adversary's whole interaction is a
        game the analysis can control.
      \item Lower bounds proved by exactly the bad-event argument later
        formalized in \cite{BR04} --- two group elements' encodings collide.
    \end{enumerate}
\end{enumerate}

\section{1998 --- mapping the space of games}

\begin{enumerate}
  \item \method \paper{Bellare}{Practice-Oriented Provable-Security}{Bell98}
    \begin{enumerate}
      \item The manifesto for the style the rest of this list assumes: explicit
        games, explicit reductions, explicit constants. Short; read it first if
        you want the philosophy before the mathematics.
    \end{enumerate}

  \item \defn \paper{Bellare, Desai, Pointcheval \& Rogaway}{Relations Among
    Notions of Security for Public-Key Encryption Schemes}{BDPR98}
    \begin{enumerate}
      \item The $\{\mathrm{IND}, \mathrm{NM}\} \times
        \{\mathrm{CPA}, \mathrm{CCA1}, \mathrm{CCA2}\}$ lattice: every
        implication proved by a game transformation, every non-implication by a
        separating counterexample.
      \item Establishes the genre of ``relations among notions'' papers ---
        papers whose entire content is the formal relationship between games.
    \end{enumerate}

  \item \turn \paper{Canetti, Goldreich \& Halevi}{The Random Oracle
    Methodology, Revisited}{CGH04}
    \begin{enumerate}
      \item A scheme provably secure in the random-oracle game and insecure for
        \emph{every} concrete instantiation of the oracle. The first hard limit
        on what winning a ROM game buys you; made concrete for a natural scheme
        in \cite{BBP03}.
    \end{enumerate}

  \item \case \paper{Cramer \& Shoup}{A Practical Public Key Cryptosystem
    Provably Secure Against Adaptive Chosen Ciphertext Attack}{CS98}
    \begin{enumerate}
      \item First efficient IND-CCA2 scheme in the standard model; the proof is
        organized as a sequence of experiments with one explicit bad event (the
        simulator wrongly rejecting a ciphertext). Shoup's house style, six
        years before he wrote the method down in \cite{Shou04}.
    \end{enumerate}
\end{enumerate}

\section{2000 --- more users, more oracles}

\begin{enumerate}
  \item \defn \paper{Bellare, Boldyreva \& Micali}{Public-Key Encryption in a
    Multi-user Setting: Security Proofs and Improvements}{BBM00}
    \begin{enumerate}
      \item \textbf{Multi-user IND-CPA/CCA}: many keys, many challenges, one
        game. The generic hybrid costs a factor of $n \cdot q_e$; the paper
        identifies when you can avoid paying it. Sixteen years later this
        decides a standard --- see \cite{BT16}.
    \end{enumerate}

  \item \defn \paper{Bellare \& Namprempre}{Authenticated Encryption: Relations
    among Notions and Analysis of the Generic Composition Paradigm}{BN00}
    \begin{enumerate}
      \item \textbf{INT-PTXT} and \textbf{INT-CTXT} games, crossed with
        IND-CPA/CCA.
      \item The verdict on Encrypt-and-MAC / MAC-then-Encrypt /
        Encrypt-then-MAC, each settled by a game argument or a counterexample.
        A definitional paper that changed what implementations do.
    \end{enumerate}

  \item \method \paper{Coron}{On the Exact Security of Full Domain Hash}{Cor00}
    \begin{enumerate}
      \item A sharper reduction for FDH: the security loss drops from the
        number of \emph{hash} queries $q_h$ (as in \cite{BR96}) to the number of
        \emph{signature} queries $q_s$, orders of magnitude smaller in practice.
      \item A \textbf{meta-reduction} showing the remaining $q_s$ loss cannot be
        removed by any reduction of that shape --- the technique behind every
        tightness lower bound in this list.
    \end{enumerate}

  \item \defn \paper{Bellare, Pointcheval \& Rogaway}{Authenticated Key Exchange
    Secure against Dictionary Attacks}{BPR00}
    \begin{enumerate}
      \item The \textbf{PAKE game}: \cite{BR93ke} plus a small password
        dictionary, where the adversary's advantage must be bounded by its
        number of \emph{online} guesses rather than made negligible. A game
        whose \emph{shape} had to change to match the primitive.
    \end{enumerate}
\end{enumerate}

\section{2001 --- a famous proof breaks}

\begin{enumerate}
  \item \turn \paper{Shoup}{OAEP Reconsidered}{Shou01}
    \begin{enumerate}
      \item The proof in \cite{BR94} is \emph{wrong}: plaintext awareness does
        not yield IND-CCA2 for a general one-way trapdoor function. One game hop
        did not hold.
      \item The most-cited cautionary tale in the field, and the strongest
        argument for writing proofs as checkable game sequences --- and,
        eventually, for machine-checking them \cite{CertiCrypt09}.
    \end{enumerate}

  \item \case \paper{Fujisaki, Okamoto, Pointcheval \& Stern}{RSA-OAEP is Secure
    under the RSA Assumption}{FOPS04}
    \begin{enumerate}
      \item The repair: a valid but much lossier chain, exploiting RSA's
        specific structure rather than generic one-wayness.
    \end{enumerate}

  \item \defn\case \paper{Boneh \& Franklin}{Identity-Based Encryption from the
    Weil Pairing}{BF03}
    \begin{enumerate}
      \item \textbf{IND-ID-CCA}: the IND-CCA game extended with an adaptive
        \emph{key-extraction} oracle for identities other than the challenge
        one. A clean example of adding one oracle to an existing game to define
        a new primitive.
      \item The proof is a textbook ROM game chain, and it is what made a whole
        decade of pairing-based definitions follow the same recipe.
    \end{enumerate}

  \item \defn\method \paper{Canetti \& Krawczyk}{Analysis of Key-Exchange
    Protocols and Their Use for Building Secure Channels}{CK01}
    \begin{enumerate}
      \item The \textbf{CK model}: a game-based key-exchange definition
        packaged with a modular composition theorem (authenticators / the
        AM-to-UM compiler).
      \item The clearest early attempt to get composition --- UC's selling
        point --- while staying inside a game.
    \end{enumerate}

  \item \paper{Canetti}{Universally Composable Security: A New Paradigm for
    Cryptographic Protocols}{Can01}
    \begin{enumerate}
      \item Listed only as the pointer to the other document: from here on the
        simulation branch is tracked in \texttt{all-uc-papers.pdf}. Much of the
        game-side work below is a response to what UC offers and games do not
        --- composition.
    \end{enumerate}
\end{enumerate}

\section{2002--2003 --- new primitives get formal games}

\begin{enumerate}
  \item \method \paper{Maurer}{Indistinguishability of Random Systems}{Mau02}
    \begin{enumerate}
      \item \textbf{Random systems} and \textbf{monotone conditions}: an
        abstract, code-free alternative to game playing, where
        ``identical-until-bad'' becomes conditional equivalence of two systems.
      \item Reaches the same bounds as code-based games by a different route;
        worth reading precisely because it is the road not taken.
        See also \cite{MPR07}.
    \end{enumerate}

  \item \defn \paper{Rogaway}{Authenticated-Encryption with Associated-Data}
    {Rog02}
    \begin{enumerate}
      \item The \textbf{AEAD game}, nonce-based: privacy and authenticity of a
        message plus unencrypted-but-authenticated headers. The game essentially
        every deployed cipher suite is now proved against.
    \end{enumerate}

  \item \method\case \paper{Cramer \& Shoup}{Design and Analysis of Practical
    Public-Key Encryption Schemes Secure against Adaptive Chosen Ciphertext
    Attack}{CS03}
    \begin{enumerate}
      \item Hash proof systems, and a long proof written as an explicit,
        numbered sequence of games --- the immediate precursor to \cite{Shou04}.
    \end{enumerate}

  \item \defn \paper{Bellare, Micciancio \& Warinschi}{Foundations of Group
    Signatures: Formal Definitions, Simplified Requirements, and a Construction
    Based on General Assumptions}{BMW03}
    \begin{enumerate}
      \item \textbf{Full-anonymity} and \textbf{full-traceability}: two games
        that provably subsume the dozen informal requirements the group
        signature literature had accumulated.
      \item The clearest demonstration of what formalizing as games buys you:
        a folklore list of properties collapses to two definitions, and the
        collapse is a theorem. Compare \cite{BSZ05} for the dynamic case.
    \end{enumerate}

  \item \defn \paper{Bellare, Namprempre, Pointcheval \& Semanko}{The
    One-More-RSA-Inversion Problems and the Security of Chaum's Blind Signature
    Scheme}{BNPS03}
    \begin{enumerate}
      \item \textbf{One-more-RSA-inversion}: a new \emph{assumption} stated as
        a game (invert $n+1$ targets given $n$ oracle inversions), because no
        standard assumption could support the proof.
      \item \textbf{Blindness} and \textbf{one-more-unforgeability} for blind
        signatures. Shows the move of pushing the interactive difficulty out of
        the proof and into a formally-stated game assumption --- now routine.
    \end{enumerate}
\end{enumerate}

\section{2004 --- the year games were formalized}

\begin{enumerate}
  \item \method \paper{Shoup}{Sequences of Games: A Tool for Taming Complexity
    in Security Proofs}{Shou04}
    \begin{enumerate}
      \item \textbf{Transitions based on indistinguishability}: replace a
        distribution with a computationally close one; the hop costs one
        assumption.
      \item \textbf{Transitions based on failure events}, justified by the
        \textbf{Difference Lemma}: if $G_i$ and $G_{i+1}$ proceed identically
        unless event $F$ occurs, then
        $\lvert \Pr[S_i] - \Pr[S_{i+1}] \rvert \le \Pr[F]$.
      \item \textbf{Bridging steps}: purely syntactic rewrites that change
        nothing but make the next real step visible. Naming these is half the
        contribution --- they are where informal proofs hide their mistakes.
      \item Worked end-to-end examples: hashed ElGamal, Cramer--Shoup, PSS.
    \end{enumerate}

  \item \method \paper{Bellare \& Rogaway}{Code-Based Game-Playing Proofs and
    the Security of Triple Encryption}{BR04,BR06}
    \begin{enumerate}
      \item \textbf{Games as pseudocode}: a game is a program with
        \texttt{Initialize}, oracle procedures and \texttt{Finalize}; the
        adversary is linked against it. Proof steps become \emph{program}
        rewrites, which is what later makes mechanization possible.
      \item The \textbf{bad flag} discipline: games are written to be
        syntactically \textbf{identical-until-bad}.
      \item The \textbf{Fundamental Lemma of game playing}: if $G$ and $H$ are
        identical-until-bad, then
        $\lvert \Pr[G^A \Rightarrow 1] - \Pr[H^A \Rightarrow 1] \rvert
        \le \Pr[G^A \text{ sets } \mathsf{bad}]$.
      \item \textbf{Lazy vs.\ eager sampling}, coin-fixing and resampling: the
        standard library of bridging rewrites, each proved once.
      \item Application: triple encryption's $\approx 2^{78}$ security bound.
      \item This and \cite{Shou04} are the two documents the phrase
        ``game-based proof'' refers to. Read both: Shoup's is the prose method,
        this one is the formal system.
    \end{enumerate}

  \item \defn\method \paper{Maurer, Renner \& Holenstein}{Indifferentiability,
    Impossibility Results on Reductions, and Applications to the Random Oracle
    Methodology}{MRH04}
    \begin{enumerate}
      \item The \textbf{indifferentiability} game, plus the composition theorem
        that ordinary game-based definitions lack: a construction
        indifferentiable from a random oracle can replace one in any
        single-stage game.
      \item The standard tool for hash-function domain extension
        (Merkle--Damg\aa{}rd, sponges) --- and, later, a cautionary tale of its
        own \cite{RSS11}.
    \end{enumerate}

  \item \defn \paper{Rogaway \& Shrimpton}{Cryptographic Hash-Function Basics}
    {RS04}
    \begin{enumerate}
      \item Seven games for hash functions --- \textbf{Coll}, \textbf{Pre},
        \textbf{ePre}, \textbf{aPre}, \textbf{Sec}, \textbf{eSec},
        \textbf{aSec} --- with the complete implication/separation map between
        them.
      \item A model of definitional hygiene: the differences between these games
        are exactly the differences practitioners get wrong.
    \end{enumerate}

  \item \defn \paper{Rogaway}{Nonce-Based Symmetric Encryption}{Rog04}
    \begin{enumerate}
      \item Replaces the random IV of \cite{BDJR97} with a \textbf{nonce} the
        adversary picks (subject to non-repetition), which is what real APIs
        expose.
      \item Read as a pair with \cite{BDJR97}: the same primitive, re-formalized
        a second time because the first game's interface did not match how the
        primitive is used.
    \end{enumerate}
\end{enumerate}

\section{2005--2007 --- mechanize it, and keep defining}

\begin{enumerate}
  \item \tool \paper{Halevi}{A Plausible Approach to Computer-Aided
    Cryptographic Proofs}{Hal05}
    \begin{enumerate}
      \item The founding document of computer-aided game-based cryptography:
        since a game sequence is a sequence of \emph{program} rewrites, a tool
        should check them.
      \item Proposes exactly the architecture --- a game-transformation engine
        plus a proof checker --- that \cite{CryptoVerif06},
        \cite{CertiCrypt09} and \cite{EasyCrypt11} went on to build.
    \end{enumerate}

  \item \tool \paper{Blanchet}{A Computationally Sound Mechanized Prover for
    Security Protocols}{CryptoVerif06}
    \begin{enumerate}
      \item First tool to \emph{find} game sequences automatically in the
        computational model (as opposed to the symbolic/Dolev--Yao model).
      \item Games as processes; hops as observational-equivalence-preserving
        process rewrites.
    \end{enumerate}

  \item \defn \paper{Rogaway \& Shrimpton}{A Provable-Security Treatment of the
    Key-Wrap Problem}{RS06}
    \begin{enumerate}
      \item \textbf{Deterministic authenticated encryption (DAE)}: a game for a
        primitive that had shipped for years (key wrapping) with no definition
        at all.
      \item The cleanest case study in reverse-engineering a formal game from
        deployed practice, rather than designing practice to fit a game.
    \end{enumerate}

  \item \defn\turn \paper{Rogaway}{Formalizing Human Ignorance:
    Collision-Resistant Hashing without the Keys}{Rog06}
    \begin{enumerate}
      \item Why there is no sensible collision-resistance \emph{game} for a
        fixed, unkeyed hash like SHA-256 --- a collision exists, so some
        adversary has it hardwired.
      \item The \textbf{explicit-reduction} workaround: state theorems as
        constructive transformations of adversaries rather than as advantage
        inequalities. A rare case of the game formalism being replaced rather
        than patched.
    \end{enumerate}

  \item \method \paper{Bellare \& Neven}{Multi-Signatures in the Plain
    Public-Key Model and a General Forking Lemma}{BN06}
    \begin{enumerate}
      \item The \textbf{general forking lemma}: \cite{PS96}'s forking stated as
        a self-contained probabilistic lemma about algorithms, decoupled from
        signatures, so it can be dropped into any game chain.
    \end{enumerate}

  \item \defn \paper{Bellare, Boldyreva \& O'Neill}{Deterministic and
    Efficiently Searchable Encryption}{BBO07}
    \begin{enumerate}
      \item \textbf{PRIV}: the best achievable privacy game once encryption must
        be deterministic --- security only for messages of high min-entropy,
        which requires the game to quantify over message distributions rather
        than message pairs.
      \item Shows how to write a formal game for a primitive that cannot meet
        the standard one, instead of declaring it insecure and moving on. (It
        is also one of the multi-stage games \cite{RSS11} later singles out.)
    \end{enumerate}

  \item \method \paper{Maurer, Pietrzak \& Renner}{Indistinguishability
    Amplification}{MPR07}
    \begin{enumerate}
      \item Composing several ``somewhat indistinguishable'' systems into a very
        indistinguishable one: the amplification counterpart to the Fundamental
        Lemma, in \cite{Mau02}'s random-systems framework.
    \end{enumerate}
\end{enumerate}

\section{2009--2012 --- machine checking, and definitional hygiene}

\begin{enumerate}
  \item \tool \paper{Barthe, Grégoire \& Zanella Béguelin}{Formal Certification
    of Code-Based Cryptographic Proofs}{CertiCrypt09}
    \begin{enumerate}
      \item \cite{BR04}'s code-based games formalized in Coq: a game is a
        program in a probabilistic imperative language, a hop is a certified
        program transformation.
      \item \textbf{Probabilistic relational Hoare logic (pRHL)}, the logic in
        which ``these two games behave identically unless bad'' becomes a
        checkable judgment. Still the technical core of EasyCrypt today.
      \item First fully machine-checked proofs of OAEP and of the Fundamental
        Lemma itself.
    \end{enumerate}

  \item \turn\defn \paper{Bellare, Hofheinz \& Kiltz}{Subtleties in the
    Definition of IND-CCA: When and How Should Challenge-Decryption be
    Disallowed?}{BHK09}
    \begin{enumerate}
      \item Four different formalizations of ``the adversary may not decrypt the
        challenge'' that the literature had treated as interchangeable, with
        separations showing they are not.
      \item The cleanest evidence that a game's small print is load-bearing.
    \end{enumerate}

  \item \tool \paper{Barthe, Grégoire, Heraud \& Zanella Béguelin}{Computer-Aided
    Security Proofs for the Working Cryptographer}{EasyCrypt11}
    \begin{enumerate}
      \item \cite{CertiCrypt09}'s ideas behind a proof assistant a cryptographer
        can actually drive: the game sequence is written in something close to
        the paper's own pseudocode, with SMT discharging the routine steps.
      \item The tool most current game-based mechanization is built on.
    \end{enumerate}

  \item \defn \paper{Boneh, Dagdelen, Fischlin, Lehmann, Schaffner \& Zhandry}
    {Random Oracles in a Quantum World}{QROM11}
    \begin{enumerate}
      \item The \textbf{quantum random oracle model (QROM)}: the adversary
        queries the oracle in superposition, which invalidates the two moves
        every \cite{BR93rom} proof relies on --- observing queries and
        programming answers.
      \item Re-proves classical results in the new game where possible, and
        separates where not. The starting point for essentially all post-quantum
        security proofs, including \cite{HHK17}.
    \end{enumerate}

  \item \turn \paper{Ristenpart, Shacham \& Shrimpton}{Careful with Composition:
    Limitations of Indifferentiability and Universal Composability}{RSS11}
    \begin{enumerate}
      \item \textbf{Single-stage vs.\ multi-stage games}: \cite{MRH04}'s
        composition theorem only applies to the former, and several important
        games (deterministic/hedged encryption, related-key attacks, proofs of
        storage) are multi-stage.
      \item The one composition theorem the game world had turns out to be
        narrower than everyone assumed --- a direct motivation for
        \cite{SSP18}.
    \end{enumerate}

  \item \defn \paper{Bellare, Ristenpart \& Tessaro}{Multi-Instance Security and
    its Application to Password-Based Cryptography}{BRT12}
    \begin{enumerate}
      \item \textbf{$m$-out-of-$n$ multi-instance games}: the adversary must
        break \emph{all} $m$ instances, which is the right question for password
        hashing and salting, and is not implied by single-instance security.
    \end{enumerate}

  \item \defn \paper{Bellare, Hoang \& Rogaway}{Foundations of Garbled Circuits}
    {BHR12}
    \begin{enumerate}
      \item Garbling promoted from a technique inside other protocols' proofs to
        a \emph{primitive} with its own games: \textbf{privacy},
        \textbf{obliviousness} and \textbf{authenticity}, in both
        indistinguishability and simulation flavors.
      \item The best example of the general move of this era: take something
        everyone proved ad hoc inside larger simulation arguments, give it
        standalone formal games, and let the rest of the literature reduce to
        them.
    \end{enumerate}

  \item \defn\case \paper{Jager, Kohlar, Schäge \& Schwenk}{On the Security of
    TLS-DHE in the Standard Model}{JKSS12}
    \begin{enumerate}
      \item The \textbf{ACCE game}: authenticated key exchange and stateful
        authenticated encryption fused into a single game, because TLS's
        handshake and record layer cannot be separated (the finished messages
        are encrypted under the key being established).
      \item The canonical example of the game being reshaped to fit a deployed
        protocol rather than the reverse.
    \end{enumerate}
\end{enumerate}

\section{2013--2016 --- new definitional idioms}

\begin{enumerate}
  \item \case \paper{Krawczyk, Paterson \& Wee}{On the Security of the TLS
    Protocol: A Systematic Analysis}{KPW13}
    \begin{enumerate}
      \item Constrained-CCA / ACCE-style treatment covering TLS's actual
        ciphersuites, including RSA key transport.
    \end{enumerate}

  \item \defn \paper{Bellare, Hoang \& Keelveedhi}{Instantiating Random Oracles
    via UCEs}{BHK13}
    \begin{enumerate}
      \item \textbf{UCE (universal computational extractor)}: a family of
        standard-model games for hash functions, strong enough to replace the
        random oracle in many proofs, defined via a two-stage
        source/distinguisher game.
      \item The most ambitious attempt to answer \cite{CGH04} with a definition
        rather than a disclaimer --- and a good illustration of how delicate a
        new game is: several UCE variants were shown too strong soon after.
    \end{enumerate}

  \item \defn\case \paper{Fischlin \& Günther}{Multi-Stage Key Exchange and the
    Case of Google's QUIC Protocol}{FG14}
    \begin{enumerate}
      \item The \textbf{multi-stage key-exchange game}: several keys per
        session, established at different times with different security levels,
        each testable. The model TLS 1.3 is analyzed in \cite{DFGS21}.
    \end{enumerate}

  \item \tool \paper{Petcher \& Morrisett}{The Foundational Cryptography
    Framework}{FCF15}
    \begin{enumerate}
      \item Another Coq embedding of code-based games, aimed at extracting
        verified implementations alongside the proof.
    \end{enumerate}

  \item \turn \paper{Bader, Jager, Li \& Schäge}{On the Impossibility of Tight
    Cryptographic Reductions}{BJLS16}
    \begin{enumerate}
      \item A general \textbf{meta-reduction} framework, in \cite{Cor00}'s
        lineage: for a broad class of schemes with unique signatures/keys,
        \emph{no} tight reduction from a non-interactive assumption exists.
      \item Turns tightness from a proof-quality question into a theorem about
        the game itself.
    \end{enumerate}

  \item \defn\case \paper{Bellare \& Tackmann}{The Multi-User Security of
    Authenticated Encryption: AES-GCM in TLS 1.3}{BT16}
    \begin{enumerate}
      \item \textbf{Multi-user AE} with concrete bounds, applied to the exact
        nonce-randomization construction TLS 1.3 was choosing at the time.
      \item \cite{BBM00}'s multi-user idea, sixteen years later, deciding a
        standard --- a good closing argument for why the constants in a game
        matter.
    \end{enumerate}

  \item \paper{Lindell}{How To Simulate It --- A Tutorial on the Simulation
    Proof Technique}{Lin16}
    \begin{enumerate}
      \item Not game-based --- included as the companion tutorial for the other
        branch, and the fastest way to see which problems each style suits.
    \end{enumerate}
\end{enumerate}

\section{2017--2019 --- modularity}

\begin{enumerate}
  \item \defn\method \paper{Hofheinz, Hövelmanns \& Kiltz}{A Modular Analysis of
    the Fujisaki--Okamoto Transformation}{HHK17}
    \begin{enumerate}
      \item The FO transform decomposed into small steps with a named game
        between each: IND-CPA $\to$ OW-PCVA $\to$ IND-CCA KEM, each hop proved
        separately and reusable.
      \item The modern standard for how to write a reduction: modular, with
        explicit intermediate notions. Every NIST post-quantum KEM's proof is
        built out of these pieces.
    \end{enumerate}

  \item \tool \paper{Basin, Lochbihler \& Sefidgar}{CryptHOL: Game-based Proofs
    in Higher-order Logic}{CryptHOL17}
    \begin{enumerate}
      \item Games as terms in Isabelle/HOL's probability monad, so game hops are
        equational rewriting rather than a bespoke program logic.
    \end{enumerate}

  \item \method \paper{Brzuska, Delignat-Lavaud, Fournet, Kohbrok \& Kohlweiss}
    {State Separation for Code-Based Game-Playing Proofs}{SSP18}
    \begin{enumerate}
      \item \textbf{Packages}: games split into stateful components with
        disjoint state and typed interfaces, composed sequentially ($\circ$) and
        in parallel ($\otimes$).
      \item \textbf{Interchange / composition lemmas} letting a sub-package be
        replaced by an indistinguishable one \emph{inside} a larger game ---
        the modularity that UC had and plain game-based proofs did not.
      \item The main structural advance in game-based proving since
        \cite{BR04}, and what makes proofs of protocol-scale objects
        (\cite{BDEFKK22}) tractable.
    \end{enumerate}

  \item \turn \paper{Drijvers, Edalatnejad, Ford, Kiltz, Loss, Neven \&
    Stepanovs}{On the Security of Two-Round Multi-Signatures}{DEFKLNS19}
    \begin{enumerate}
      \item Several published two-round Schnorr multi-signature proofs are
        broken, by a concurrent (Wagner/ROS) attack their games' rewinding
        arguments could not see.
      \item The modern counterpart to \cite{Shou01}: the reason the
        multi-signature literature now insists on fully formal games. See
        \cite{NRS21} for the repair.
    \end{enumerate}

  \item \tool \paper{Canetti, Stoughton \& Varia}{EasyUC: Using EasyCrypt to
    Mechanize Proofs of Universally Composable Security}{EasyUC19}
    \begin{enumerate}
      \item The two branches meet in one tool: UC simulation proofs carried out
        in the game-based prover. Useful for seeing exactly where the styles
        differ formally.
    \end{enumerate}
\end{enumerate}

\section{2019--2021 --- tooling consolidates}

\begin{enumerate}
  \item \tool \paper{Barbosa, Barthe, Bhargavan, Blanchet, Cremers, Liao \&
    Parno}{SoK: Computer-Aided Cryptography}{SoKCAC21}
    \begin{enumerate}
      \item The survey to read before picking a tool: design-level (game-based)
        provers, symbolic provers, and verified implementations, with what each
        can and cannot do.
    \end{enumerate}

  \item \tool \paper{Haselwarter et al.}{SSProve: A Foundational Framework for
    Modular Cryptographic Proofs in Coq}{SSProve21}
    \begin{enumerate}
      \item \cite{SSP18}'s state-separating proofs, mechanized: packages and
        their composition laws as a Coq library with a relational program logic
        underneath.
      \item Closes the loop \cite{Hal05} opened in 2005 --- modular \emph{and}
        machine-checked.
    \end{enumerate}

  \item \tool \paper{Baelde, Delaune, Jacomme, Koutsos \& Moreau}{An Interactive
    Prover for Protocol Verification in the Computational Model}{Squirrel21}
    \begin{enumerate}
      \item A different bet: a logic for computational indistinguishability
        (after Bana--Comon) instead of explicit game rewriting.
    \end{enumerate}

  \item \defn\case \paper{Nick, Ruffing \& Seurin}{MuSig2: Simple Two-Round
    Schnorr Multi-Signatures}{NRS21}
    \begin{enumerate}
      \item A two-round scheme with a full game-based proof that survives the
        attack of \cite{DEFKLNS19}, in the ROM (plus a tighter argument in the
        algebraic group model).
      \item Worth reading immediately after \cite{DEFKLNS19}: the same game, one
        broken proof and one repaired one, a decade of technique apart.
    \end{enumerate}

  \item \case \paper{Dowling, Fischlin, Günther \& Stebila}{A Cryptographic
    Analysis of the TLS 1.3 Handshake Protocol}{DFGS21}
    \begin{enumerate}
      \item \cite{FG14}'s full multi-stage game applied to the shipped TLS 1.3
        handshake --- and an example of a game-based analysis that shaped a
        standard while it was being written.
    \end{enumerate}
\end{enumerate}

\section{2022 and after --- games at protocol scale}

\begin{enumerate}
  \item \case\method \paper{Brzuska, Delignat-Lavaud, Egger, Fournet, Kohbrok \&
    Kohlweiss}{Key-schedule Security for the TLS 1.3 Standard}{BDEFKK22}
    \begin{enumerate}
      \item \cite{SSP18} applied at full scale: the TLS 1.3 key schedule proved
        as a composition of packages, at a size no monolithic game sequence
        would survive.
      \item The current answer to ``can game-based proofs handle real
        protocols?'' --- yes, but only modularly.
    \end{enumerate}
\end{enumerate}

\section{Appendix: the standing critiques}

\noindent
Not definitional papers, but the arguments about what a game-based theorem
actually claims. Useful once the rest of the list makes sense.

\begin{enumerate}
  \item \paper{Koblitz \& Menezes}{Another Look at ``Provable Security''}{KM04}
    Non-tight reductions read as meaningful, and the gap between a model's
    oracles and reality.
  \item \paper{Bellare, Boldyreva \& Palacio}{An Uninstantiable
    Random-Oracle-Model Scheme for a Hybrid Encryption Problem}{BBP03}
    \cite{CGH04}'s separation, for a natural scheme rather than a contrived one.
  \item \paper{Dent}{Fundamental problems in provable security and
    cryptography}{Dent06} A survey of what the enterprise cannot deliver,
    written from inside it.
  \item \paper{Bernstein \& Lange}{Non-uniform Cracks in the Concrete: The Power
    of Free Precomputation}{BL13} Adversaries with free preprocessing beat many
    concrete bounds as stated, so the quantifiers in a game matter as much as
    its shape.
\end{enumerate}

\vspace{1em}
\noindent\rule{\linewidth}{0.4pt}

\medskip\noindent
\textbf{The list in one line.} Games were how definitions got written from 1982
\cite{GM82}, and how primitives got \emph{fully} formalized from 1997
\cite{BDJR97}, but they were not themselves an object of study until 2004, when
\cite{Shou04} and \cite{BR04} gave the sequence-of-games proof a grammar. What
follows splits in two: mechanizing that grammar (\cite{Hal05} $\to$
\cite{CryptoVerif06} $\to$ \cite{CertiCrypt09} $\to$ \cite{EasyCrypt11} $\to$
\cite{CryptHOL17,SSProve21}), and buying back the composability the style gave
up when it declined to use simulators (\cite{CK01} $\to$ \cite{MRH04} $\to$ its
limits in \cite{RSS11} $\to$ state separation in \cite{SSP18} $\to$
\cite{BDEFKK22}).

\bibliographystyle{alpha}
\bibliography{refs}

\end{document}
